\section*{अध्याय \ref{ch:g1:healthy-habits} --- सेहतमंद आदतें}

\begin{solution}{exo:g1:healthy-habits:1}
\omterm{def:g1:healthy-habits:germs}{रोगाणु} इतने छोटे \omterm{def:g1:living-or-not:living}{सजीव} हैं कि दिखाई ही नहीं देते, और वे हर जगह मिलते
हैं; कुछ शरीर के भीतर पहुँचकर हमें बीमार कर सकते हैं। दिखाई न देने के
कारण हाथ बिलकुल साफ़ दिख सकते हैं और फिर भी उन्हें लिए हो सकते हैं।
\end{solution}

\begin{solution}{exo:g1:healthy-habits:2}
खाने से पहले, शौचालय के बाद, और बाहर खेलने या \omterm{def:g1:animals-around-us:animal}{जंतुओं} को छूने के बाद।
\end{solution}

\begin{solution}{exo:g1:healthy-habits:3}
धीरे-धीरे $20$ तक गिनने जितनी देर (लगभग एक छोटा गीत)। आसानी से भूली
जाने वाली जगहें: उँगलियों के बीच, अँगूठों के चारों ओर, और हाथों की
ऊपरी सतह।
\end{solution}

\begin{solution}{exo:g1:healthy-habits:4}
खाने के बाद दाँतों पर भोजन के कण रह जाते हैं। मुँह के \omterm{def:g1:healthy-habits:germs}{रोगाणु} इन कणों
पर पलते हैं, ख़ास तौर पर मीठे कणों पर, और थोड़ा-थोड़ा करके दाँतों में
छेद बना देते हैं, जो दुखते हैं। ब्रश करना इन कणों को दावत से पहले ही
बुहार देता है।
\end{solution}

\begin{solution}{exo:g1:healthy-habits:5}
नाश्ते के बाद और सोने से पहले, हर रोज़ --- और इसमें लगभग एक गीत जितना
समय लगना चाहिए, हर दाँत की हर सतह पर।
\end{solution}

\begin{solution}{exo:g1:healthy-habits:6}
\omterm{prop:g1:healthy-habits:needs}{विविध भोजन} (नाश्ते में फल, पीने को पानी), हलचल (बाहर खेलना, साइकिल
चलाना), और \omterm{prop:g1:healthy-habits:needs}{नींद} (जल्दी सोकर एक लंबी रात)। उदाहरण खुले हैं।
\end{solution}

\begin{solution}{exo:g1:healthy-habits:7}
वह आराम करता है, अपनी मरम्मत करता है और बढ़ता है। लंबी रात बच्चे को
चाहिए --- बड़ों से कहीं लंबी।
\end{solution}

\begin{solution}{exo:g1:healthy-habits:8}
जैसे: जम्हाइयाँ, जलती या भारी \omterm{def:g1:the-five-senses:sense}{आँखें}, बेवजह चिड़चिड़ाहट।
\end{solution}

\begin{solution}{exo:g1:healthy-habits:9}
दूध के दाँतों को गिरने से पहले कई साल चबाना पड़ता है; उनमें छेद बड़ों
वाले दाँतों जितना ही दुखता है; और नीचे इंतज़ार कर रहे बड़ों वाले दाँत
साफ़, सेहतमंद मुँह में आने के हक़दार हैं। इनमें से दो कारण काफ़ी हैं।
\end{solution}

\begin{solution}{exo:g1:healthy-habits:10}
चमकी \omterm{def:g1:healthy-habits:germs}{रोगाणुओं} को दर्शाती है। हाथ मिलाना, दरवाज़े खोलना और पेंसिलें
साझा करना उन्हें बिना दिखे एक हाथ से दूसरे हाथ तक पहुँचा देता है ---
और $20$ की धीमी गिनती तक साबुन से धोना उन्हें हटा देता है, चमकी और
\omterm{def:g1:healthy-habits:germs}{रोगाणु} दोनों को।
\end{solution}
